\chapter{Proofs for Full Natural-Deduction Normalization}
\label{annex:p07-natural-deduction-normalization}

\begin{annexitem}{annex:p7.nj.nonvacuous-obstruction}{thm:p07-nj-nonvacuous-obstruction}{A closed derivation with a blocked sole detour}
\begin{proof}
Start with two distinct leaves \([A]^i,[B]^j\), where
\(A,B\) are distinct atoms. Conjunction introduction gives
\(A\land B\), and its first projection gives \(A\).
The \(B\)-leaf is still an open occurrence above this conclusion:
projection has not discharged it. Implication introduction with
label \(j\) therefore legally derives \(B\to A\), discharging
that occurrence. A second implication introduction with label
\(i\) gives the closed conclusion \(A\to(B\to A)\).
In proof-term notation the tree is
\[
       \lambda u^A.\lambda v^B.\pi_1\langle u,v\rangle.
\]
Both abstractions are nonvacuous in this unreduced term.

The only introduction conclusion used as a matching elimination
major premise is the pair \(\langle u,v\rangle\) under
\(\pi_1\). Its principal contraction returns \(u\).
In the surrounding tree the result would be
\(\lambda u^A.\lambda v^B.u\), whose inner implication
introduction discharges no \(B\)-occurrence. It is not a
\(\mathsf{NJ}_{\mathrm{nv}}\) derivation. There are no other
eliminations, disjunction cases or existential cases in the
tree, hence no other principal or permitted permutation step.
Alpha-renaming changes none of these occurrence counts.
Thus there is no legal first reduction on a route from this
tree to a detour-free tree within the restricted calculus.

In \(\mathsf{NJ}_{\mathrm v}\), the empty inner discharge
is legal. The same contraction reaches
\(\lambda u^A.\lambda v^B.u\), a closed derivation with no
introduction--elimination or permutation detour. This proves
the asserted difference without identifying the two calculi.
\end{proof}
\end{annexitem}

\begin{annexitem}{annex:p7.nj.consequence-bridge}{thm:p07-nj-vacuous-consequence-bridge}{Simulating every empty discharge by a genuine occurrence}
\begin{proof}
Every rule of the nonvacuous presentation is already a permitted
instance of the empty-discharge presentation, so one direction
leaves the tree unchanged.

For the other direction use induction on the finite source tree.
Convert its premise derivations first. No conversion in this
induction removes an original open leaf, and every new leaf
described below will be discharged by the rule currently being
translated. Hence original nonempty discharge classes remain
nonempty. A nondischarging rule, or a rule whose required classes
already are nonempty, can be copied with its original conditions.

For \(\to I\) with an empty class of \(A\)-assumptions,
let the converted premise derive \(B\). Add a fresh leaf
\([A]^i\), apply \(\land I\) to obtain \(B\land A\),
then \(\land E_1\) to return to \(B\). The new assumption
occurrence is open above this premise and is discharged by
\(\to I_i\). The desired conclusion \(A\to B\) and all
original open occurrences are retained.

For \(\lor E\), perform that same construction separately
in each minor derivation whose discharged class is empty.
In the \(A\)-branch insert a new \([A]^i\), combine the
branch conclusion \(C\) with it to get \(C\land A\),
and project back to \(C\). In the \(B\)-branch use a new
\([B]^j\) and \(C\land B\). The existing \(\lor E\)
now discharges a nonempty class in each branch. Its major
premise is unchanged. If just one branch was empty, only
that branch needs the construction.

For \(\exists E\), let the major formula be \(\exists x A\)
and the eigenvariable \(y\). If its minor derivation of
\(C\) has no discharged \(A[y/x]\)-occurrence, append a
fresh \([A[y/x]]^i\), form \(C\land A[y/x]\), and project
back to \(C\). The final \(\exists E_i\) discharges this
new occurrence. Its eigenvariable remains absent from the
conclusion, major premise formula and all other open assumption
formulas, exactly as before. The new occurrence is the one
explicitly permitted by that rule's eigenvariable condition.

These added leaves are introduced below the internal rule
applications of the converted minor derivations. They do not
become open assumptions above an earlier \(\forall I\)
or \(\exists E\), so they cannot invalidate its freshness.
The newly added steps are only conjunction introduction and
projection and impose no new eigenvariable requirement.
Fresh discharge labels avoid capture of any old class.
Every added leaf is removed from the open context by its
specified final discharge; none remains open at the root.
All original open occurrences survive once each. Each source
node needs finitely many additions, so the induction produces
a finite target tree. Its inserted projection detours also
show why this consequence translation does not preserve normality.
\end{proof}
\end{annexitem}

\begin{annexitem}{annex:p7.nj.substitution}{thm:p07-nj-substitution}{Proof and individual substitution for all constructors}
\begin{proof}
Finite trees use only finitely many names. Rename bound proof
labels and local individual eigenvariables to avoid the finite
sets of names in the terms and proofs about to be inserted.
Copies receive disjoint internal names; free proof classes are
left unchanged. This preliminary renaming respects the formula
annotations and all eigenvariable conditions.

For proof substitution, induct on the derivation of \(M\).
A selected assumption leaf is replaced by a copy of \(N:A\);
an unselected leaf is unchanged. Each conjunction introduction,
projection, injection, implication application and absurdity
elimination is reconstructed from its transformed premises.
Their premise and conclusion formulas are unchanged. At
\(\to I\), avoid capture of free labels in \(N\) by its
bound label, substitute in the body and perform the same
discharge. At \(\lor E\), substitute in the major and both
minor derivations, treating their two bound classes locally,
then reconstruct the rule. A class that becomes empty remains
permitted in \(\mathsf{NJ}_{\mathrm v}\).

At \(\forall I\), its renamed eigenvariable is fresh for
the inserted proof and for every assumption that can be added
by that proof. It remains absent from all open assumptions of
the new premise. Universal elimination and existential
introduction retain their same witness terms and formula
instances, so their transformed premises still have the
required formulas. At \(\exists E\), the renamed eigenvariable
is absent from the inserted proof, the resulting conclusion,
the major formula and all external open assumptions. Its
one bound proof class is kept separate from the substitution
domain. Reconstructing the rule therefore preserves every
freshness and discharge condition. This exhausts the rules.
Only copies of the given \(N\) can contribute additional
open leaves, yielding the asserted containment. If the selected
class has no occurrences, no copy is inserted.

For individual substitution use a separate induction on the
derivation tree. Terms and formulas use the fixed total
capture-avoiding operation. An assumption becomes the
corresponding substituted assumption, and all Boolean
introduction/elimination rules keep their shape with every
formula uniformly substituted. The same labeled occurrences
are discharged. At \(\forall I\), let the eigenvariable
be \(z\), chosen different from the substituted variable
\(x\) and absent from the substituting term \(t\).
Uniform substitution cannot introduce \(z\) into an open
assumption, and
\((\forall z B)[t/x]\) has representative
\(\forall z(B[t/x])\). Thus universal introduction is legal.
At \(\forall E\), substitution in the instantiating term
and formula commutes with the original instantiation, after
fresh bound-variable renaming. Both sides are the same total
capture-avoiding simultaneous substitution on the original
formula. The existential-introduction case uses this same
substitution identity.

At \(\exists E\), choose its eigenvariable \(z\) distinct
from \(x\) and absent from \(t\). It remains absent from
the substituted conclusion, major formula and external
assumptions. The discharged class becomes precisely the
corresponding substituted matrix instance, so the existential
elimination rule can be reconstructed. These arguments also
handle a vacuous individual quantifier and any zero-occurrence
discharge class. Bound-variable alpha-equivalence resolves
the choice of fresh representatives without changing any free
variable of the asserted conclusion or assumptions.
\end{proof}
\end{annexitem}

\begin{annexitem}{annex:p7.nj.principal-preservation}{thm:p07-nj-principal-preservation}{Every principal conversion and its surrounding proof context}
\begin{proof}
A projection of an introduced pair retains the selected
component derivation and deletes the other. Its formula is
the elimination conclusion, and its open assumptions are a
subset of the original union. An introduced implication
followed by application substitutes the minor proof of \(A\)
for the bound \(A\)-class in the major proof's body. The
proof-substitution lemma gives the conclusion \(B\) with
exactly the allowed source assumption formulas.

For a case analysis on a left injection, retain its left
minor derivation and substitute the injection's proof of
\(A\) for the class \(u:A\). The other branch is deleted.
The right injection uses the right branch and its class
\(v:B\). The substitution lemma gives \(C\) in both
cases, including when the chosen branch's class was empty.

A universally introduced proof followed by instantiation at
\(t\) gives the body proof with \(t\) substituted for its
universal eigenvariable. That variable was free in no open
assumption, so the individual-substitution lemma gives the
required instance without changing the external assumptions.
For an existential package followed by unpacking, first
replace the unpacking eigenvariable \(y\) by its actual
witness \(t\) in the minor proof. Freshness makes this a
legal individual substitution and leaves its conclusion and
external assumptions unchanged. Its discharged assumption
now has formula \(A[t/x]\), exactly the package proof's
formula. Substitute that proof for class \(u\). The two
substitution lemmas give the original conclusion \(C\),
and no assumption other than those already in the package
and external minor proof appears.

To lift a local conversion through a proof context, reconstruct
the rule immediately below its occurrence and proceed towards
the root. Each conclusion formula at this boundary is unchanged.
The transformed subtree can only lose source open assumption
formulas; copied assumptions retain their preexisting surrounding
discharge ownership after internal fresh renaming. A later
discharge may therefore become empty, which is allowed.
An eigenvariable that was absent from the original external
open assumptions remains absent from their subset. Internal
names of moved or copied proofs were made fresh before the
conversion, so no capture is introduced. Every surrounding
rule remains legal. This proves preservation at every proof-term
position, not just at the root of a derivation.
\end{proof}
\end{annexitem}

\begin{annexitem}{annex:p7.nj.permutation-preservation}{thm:p07-nj-permutation-preservation}{Every following elimination, with all eigenvariable cases}
\begin{proof}
In the case permutation, both original minor derivations end at
the same formula \(C\), and the following frame turns a proof
of \(C\) into a proof of \(B\). Apply a fresh copy of this
frame to each minor conclusion, and reapply the original case
rule with new conclusion \(B\). In the unpacking permutation,
apply the following frame to its one minor conclusion and
reapply the original unpacking rule. We check all seven frames.
\begin{enumerate}
\item For \(\pi_1\), the common formula is \(C_1\land C_2\)
and each new branch concludes \(C_1\). No auxiliary derivation
or new binder is present.
\item For \(\pi_2\), the same argument gives \(C_2\).
These two checks establish four concrete permutations in total.
\item For implication elimination, the common formula is
\(C_1\to B\). Append the existing minor proof of \(C_1\)
to each relevant branch and apply \(\to E\). Its copies
have fresh internal binders and eigenvariables. The earlier
branch labels are fresh for these copies, so their formerly
external assumptions do not become accidentally discharged.
In the unpacking column its eigenvariable \(y\) is also fresh
for the copied proof's assumptions and formulas, so the final
\(\exists E\) remains legal.
\item For a following case rule, the common formula is
\(C_1\lor C_2\). Append its two minor derivations of \(B\)
to each new occurrence of that major formula. Their bound
classes \(c,d\) are first freshened away from the earlier
case or unpacking classes. When copied into two earlier case
branches they receive separate \(c_1,d_1,c_2,d_2\) names.
Each new later case discharges just its own classes; the earlier
case still discharges \(u,v\), and the earlier unpacking
still discharges \(u\). In the latter column fresh \(y\)
is absent from all new external assumptions and from \(B\).
\item For universal elimination, the common formula is
\(\forall x C_0\). Instantiate each branch at the same
outside term \(t\), obtaining \(C_0[t/x]\). An earlier
unpacking eigenvariable must first be fresh for \(t\) and
for this final instance, as well as for the surrounding data.
Free occurrences of its old spelling inside \(t\) are not
renamed. Thus moving the instantiation does not capture them.
The registered total substitution determines the same instance
in each branch.
\item For a following unpacking, the common formula is
\(\exists x C_0\). Append its minor derivation of \(B\)
under a fresh \(c:C_0[z/x]\). Its eigenvariable \(z\)
is chosen fresh for the earlier branch assumptions as well
as every other displayed formula and term. Those earlier
assumptions had already been discharged before the following
unpacking on the left, so its former side condition alone
would not have guaranteed their freshness. In the earlier-case
column choose distinct \(z_1,z_2\) and \(c_1,c_2\) in the
two copies. In the earlier-unpacking column choose \(z\ne y\)
and \(c\ne u\). The new inner rule discharges \(c\) and
the retained outer rule discharges \(u\); both eigenvariable
absence conditions now hold, with neither variable escaping
its branch.
\item For absurdity elimination, the common formula is
\(\bot\). Apply \(\bot E\) to obtain \(B\) in each
branch. This step introduces no binder. Nevertheless an earlier
unpacking eigenvariable must be fresh for \(B\), whose free
parameters are arbitrary; this is included in the prescribed
renaming before the move.
\end{enumerate}
The list covers seven following eliminations in each of the
two columns. The external assumptions of the result are
contained in the union already present in the original major,
branches and auxiliary derivations. Copying duplicates occurrences,
not new assumption formulas. The earlier labels capture exactly
their original branch occurrences; all new internal labels are
distinct. These facts prove the local preservation assertion.

For absurdity absorption, the retained proof \(M\) already
derives \(\bot\). Each of the seven frames is replaced by
the allowed \(\bot E\) to its original final formula.
For projections this is the selected conjunct; for application
the consequent; for case or unpacking their branch conclusion;
for universal elimination the same capture-avoiding instance;
and for a second absurdity elimination its final formula.
Auxiliary assumptions disappear with their discarded proofs.
An eigenvariable belonging to a discarded unpacking branch was
absent from its final conclusion and never bound any variable
of the original major proof \(M\), so its deletion causes
no escaping occurrence. This proves all seven absorption cases.

Finally reconstruct the surrounding proof context as in the
principal-preservation argument. Its premise conclusion is
unchanged, its open support can only shrink, copied free classes
keep the same external discharge ownership, and an enclosing
class that becomes empty remains licensed. Every local
preservation result therefore holds at arbitrary depth.
\end{proof}
\end{annexitem}

\begin{annexitem}{annex:p7.nj.cps-simulation}{thm:p07-nj-cps-simulation}{The complete quantified continuation simulation}
\begin{proof}
First, structural induction on annotated proof terms gives
\[
 e(M[t/x])=e(M),\qquad
 e(M[N/u])=e(M)[e(N)/u],
\]
modulo renaming bound proof variables. Individual substitution
changes only the term and formula information deliberately
erased by \(e\). At \(\Lambda\), instantiation, package and
unpacking it leaves the corresponding retained constructor;
the proof-variable binders in the second identity use the same
capture-avoiding substitutions on each side. At assumptions
the identities are their definitions, and the remaining
constructors follow from their children. Induction on formulas
also proves \(A[t/x]^{\mathrm c}=A^{\mathrm c}\) and
\(A[t/x]^{\mathrm v}=A^{\mathrm v}\): predicate arguments
are ignored and every logical constructor is retained.

\paragraph{Typing every constructor.}
Prove simultaneously that \(\mathcal T(e(M)):A^{\mathrm c}\)
and that \(\mathcal C(e(M),K):o\) whenever
\(K:A^{\mathrm v}\to o\), under the translated open context
and the marker \(\star:\iota\). The induction follows the
source derivation, so all formula types of subterms are known.
For an assumption \(u:A\), the target variable has type
\((A^{\mathrm v}\to o)\to o\), and \(uK:o\).
For implication introduction the value
\(\lambda u.\mathcal T(V)\) has type
\(B^{\mathrm c}\to C^{\mathrm c}\); its continuation accepts
exactly that type. For application, a value
\(f:B^{\mathrm c}\to C^{\mathrm c}\) accepts
\(\mathcal T(W):B^{\mathrm c}\), and its result accepts
the \(C\)-continuation \(K\), giving \(o\).

For a pair, \(\lambda p.p\mathcal T(V)\mathcal T(W)\)
has type \((B^{\mathrm c}\to C^{\mathrm c}\to o)\to o\).
The projection continuations use
\(\lambda u.\lambda v.uK\) and
\(\lambda u.\lambda v.vK\), respectively, of type
\(B^{\mathrm c}\to C^{\mathrm c}\to o\), with \(K\)
at the selected component's value type. For a left injection,
\(\lambda u.\lambda v.u\mathcal T(V)\) has type
\((B^{\mathrm c}\to o)\to(C^{\mathrm c}\to o)\to o\);
the right injection uses the second continuation. For case
analysis, its two target branches
\(\lambda u.\mathcal C(W,K)\) and
\(\lambda v.\mathcal C(Z,K)\) have the two continuation
types just displayed, so their application to the sum value
again has type \(o\).

Universal introduction yields a value
\(\lambda\xi.\mathcal T(V):\iota\to B^{\mathrm c}\).
Universal elimination applies such a value to
\(\star:\iota\), then to the instance continuation, which
has the required type because individual substitution does
not change translated types. A package yields
\(\lambda p.p\star\mathcal T(V)\) of type
\((\iota\to B^{\mathrm c}\to o)\to o\).
Unpacking supplies the continuation
\(\lambda\xi.\lambda u.\mathcal C(W,K)\) of type
\(\iota\to B^{\mathrm c}\to o\). Its dummy individual
argument changes no conclusion type. Finally, the absurdity
premise has computation type \((o\to o)\to o\).
The argument \((\lambda z.\lambda r.r)K\) has type
\(o\to o\), where \(z\) has \(K\)'s type and \(r:o\).
Thus the last row is also typed. These are all fourteen
constructor cases, counting both projections and injections.
Wrapping the \(\mathcal C\)-judgment in fresh \(\lambda k\)
establishes the simultaneous \(\mathcal T\)-judgment.

\paragraph{The continuation and substitution identities.}
The wrapper gives \(\mathcal T(U)K\to_\beta\mathcal C(U,K)\):
its bound \(k\) is fresh for \(U\), and substitution in the
continuation positions of the defining table gives exactly
the right side, up to fresh binder names.

If \(K\to_\beta L\), then
\(\mathcal C(U,K)\to_\beta^+\mathcal C(U,L)\) for every
erased \(U\). Induct on \(U\). In the variable and
introduction rows \(K\) occurs directly, so reduce there.
Application, projections and universal elimination put \(K\)
in the continuation supplied to a proper major subterm;
apply the induction hypothesis to that subterm. Case and
unpacking put copies of \(K\) into their branch continuation
calls. Apply the branch hypotheses, reducing every one of
the finitely many copies, then use the major hypothesis.
Every constructor retains at least one copy. Absurdity retains
\(K\) as the argument of \((\lambda z.\lambda r.r)K\);
reduce inside that argument before contracting the erasing
redex, and apply the major hypothesis. This explains why
the administrative redex in the definition was not simplified.
A finite nonempty reduction in \(K\) lifts by repeating this
argument; an empty reduction lifts by reflexivity.

For capture-avoiding proof substitution, the following two
statements hold by simultaneous structural induction on \(U\):
\[
\begin{aligned}
 \mathcal C(U,K)[\mathcal T(V)/u]
 &\to_\beta^*
   \mathcal C(U[V/u],K[\mathcal T(V)/u]),\\
 \mathcal T(U)[\mathcal T(V)/u]
 &\to_\beta^*\mathcal T(U[V/u]).
\end{aligned}
\]
If \(U=u\), the first left side is
\(\mathcal T(V)K[\mathcal T(V)/u]\), whose wrapper
contraction gives the assertion. A different variable is
unchanged. At an introduction, substitution passes through
fresh target binders and applications, and the inductive
\(\mathcal T\)-statements apply to every child. At application,
projection, universal elimination or absurdity, apply the
\(\mathcal C\)-induction hypothesis to its major operand
with the entire displayed continuation. If that continuation
contains a recursive \(\mathcal T\)-call, apply its
induction hypothesis too, and lift its reductions through
the major \(\mathcal C\)-call by the preceding continuation
property. At case and unpacking, apply the branch
\(\mathcal C\)-hypotheses inside the supplied continuation
and lift them in the same way, then the major hypothesis.
These checks include both branches of a case and both
propositional and quantified introductions. They do not
assume that a derived continuation is free of \(u\): the
substitution shown in its argument is essential. Wrapping
in fresh \(\lambda k\) gives the second simultaneous assertion.

\paragraph{Every root principal step.}
For implication, the translated continuation expression is
\[
 (\lambda f.f\mathcal T(V)K)(\lambda u.\mathcal T(U)).
\]
It reduces first to
\((\lambda u.\mathcal T(U))\mathcal T(V)K\), then to
\(\mathcal T(U)[\mathcal T(V)/u]K\). The substitution
identity and wrapper contraction yield
\(\mathcal C(U[V/u],K)\). At least the first beta contraction
has occurred, even if \(u\) had no occurrences.

For the first pair projection, apply its continuation
\(\lambda p.p(\lambda u.\lambda v.uK)\) to the pair value
\(\lambda q.q\mathcal T(U)\mathcal T(V)\).
Successive beta contractions give
\((\lambda u.\lambda v.uK)\mathcal T(U)\mathcal T(V)\),
then \(\mathcal T(U)K\), then \(\mathcal C(U,K)\).
The second projection uses \(\lambda u.\lambda v.vK\)
and selects \(\mathcal T(V)K\), then \(\mathcal C(V,K)\).
Both paths are positive.

For a left injection followed by case analysis, apply
\(\lambda s.s(\lambda u.\mathcal C(V,K))
                  (\lambda v.\mathcal C(W,K))\)
to \(\lambda a.\lambda b.a\mathcal T(U)\).
The beta steps select
\((\lambda u.\mathcal C(V,K))\mathcal T(U)\), then
\(\mathcal C(V,K)[\mathcal T(U)/u]\).
Substitution gives \(\mathcal C(V[U/u],K)\), taking the
bound \(u\) fresh for \(K\). A right injection selects
\((\lambda v.\mathcal C(W,K))\mathcal T(U)\) and yields
\(\mathcal C(W[U/v],K)\). These are two separate positive
simulation paths.

For universal introduction followed by instantiation,
\[
 \mathcal C(\mathsf{ae}(\mathsf{ai}(U)),K)
  =(\lambda f.f\star K)(\lambda\xi.\mathcal T(U))
  \to_\beta^+\mathcal T(U)K
  \to_\beta\mathcal C(U,K).
\]
The erased actual contractum \(M[t/x]\) is \(e(M)\), by
the erasure identity; the proof constructors, and hence the
positive simulation steps, have not disappeared.

For unpacking a package the corresponding expression is
\[
 (\lambda s.s(\lambda\xi.\lambda u.\mathcal C(V,K)))
       (\lambda p.p\star\mathcal T(U)).
\]
It reduces to
\((\lambda\xi.\lambda u.\mathcal C(V,K))\star\mathcal T(U)\),
then to \((\lambda u.\mathcal C(V,K))\mathcal T(U)\),
then to \(\mathcal C(V,K)[\mathcal T(U)/u]\), and by
substitution to \(\mathcal C(V[U/u],K)\). Individual
substitution in the actual existential contractum erases;
proof substitution commutes with erasure. This proves the
seventh principal row. No bottom-introduction case exists.

These simulations lift through every source proof context.
Induct on its one-hole structure using the full translation
table. A hole in a major operand uses the simulation hypothesis
for its particular continuation. A hole in a minor operand,
introduction operand or branch gives reductions inside one
or more displayed \(\mathcal T\) or \(\mathcal C\) calls.
Reduce the finitely many copies and use strict continuation
closure whenever they occur inside a continuation passed to
another \(\mathcal C\)-call. The absurdity case is covered
by its retained continuation argument. The resulting target
path remains nonempty in all cases. Wrapping the whole term
in its fresh \(\lambda k\) proves the principal assertion
for \(\mathcal T\), including arbitrary source depths.

\paragraph{Every root permutation.}
For an erased elimination frame write \(H_E(K)\) for the
following continuation:
\[
\begin{array}{c|l}
E[-]&H_E(K)\\ \hline
[-]L&\lambda f.f\mathcal T(L)K\\
\pi_1[-]&\lambda p.p(\lambda u.\lambda v.uK)\\
\pi_2[-]&\lambda p.p(\lambda u.\lambda v.vK)\\
\mathsf{ae}([-])&\lambda f.f\star K\\
\mathsf{case}([-];c.L;d.R)&
 \lambda s.s(\lambda c.\mathcal C(L,K))(\lambda d.\mathcal C(R,K))\\
\mathsf{unpack}([-];c.L)&
 \lambda s.s(\lambda\xi.\lambda c.\mathcal C(L,K))\\
\mathsf{abort}([-])&(\lambda z.\lambda r.r)K
\end{array}
\]
The defining table gives
\(\mathcal C(E[U],K)=\mathcal C(U,H_E(K))\) for each
of these seven frames. Substituting a case term as \(U\)
expands both sides of its permutation to
\[
 \mathcal C(M,\lambda s.s
       (\lambda u.\mathcal C(N,H_E(K)))
       (\lambda v.\mathcal C(P,H_E(K)))).
\]
Substituting an unpacking term expands both sides to
\[
 \mathcal C(M,\lambda s.s
       (\lambda\xi.\lambda u.\mathcal C(N,H_E(K)))).
\]
The fresh-name conditions of the permutation table make
these syntactic identities modulo alpha-equivalence, including
the copied minor derivations and dummy individual binders.
No beta contraction is used for this identity. The seven
frame calculations therefore cover all fourteen permutations.
Induction through the remaining source context and the
\(\mathcal T\)-wrapper preserves the identity, proving
the contextual permutation assertion.
\end{proof}
\end{annexitem}

\begin{annexitem}{annex:p7.nj.permutation-decrease}{thm:p07-nj-permutation-decrease}{A strict measure despite duplication}
\begin{proof}
All values of \(n,w\) are positive by induction on their
defining table: the variable case is one, and every remaining
row is a positive value, a sum of positive values, or a
product of positive values. Substitution of each of the
seven frame rows into the table gives
\(w(E[U])=w(U)+n(U)a_E\) and
\(n(E[U])=n(U)b_E\), with the displayed positive
coefficients. Projections, universal elimination and absurdity
have coefficients \(1,1\); application uses \(w(L),1\);
case uses \(w(L)+w(R),2(n(L)+n(R))\); and unpacking
uses \(w(L),2n(L)\). This checks every elimination row.

For a case permutation, abbreviate the sizes of its three
subproofs by their indicated \(n\)'s and \(w\)'s. Both
sides have
\[
             n=2n(M)(n(N)+n(P))b_E.
\]
The left side's \(w\) is
\[
 w(M)+n(M)(w(N)+w(P))
       +2n(M)(n(N)+n(P))a_E,
\]
and the right side's is
\[
 w(M)+n(M)(w(N)+w(P))
       +n(M)(n(N)+n(P))a_E.
\]
Subtracting gives the claimed strictly positive difference.
The values of a copied frame agree with those of the original
frame because renaming changes neither recurrence. This
calculation includes copying both later-case branches or a
later-unpacking minor proof; it makes no claim that ordinary
tree size decreases.

For an unpacking permutation both \(n\)-values are
\(2n(M)n(N)b_E\). The left \(w\)-value is
\(w(M)+n(M)w(N)+2n(M)n(N)a_E\), whereas the right is
\(w(M)+n(M)w(N)+n(M)n(N)a_E\).
Their difference is \(n(M)n(N)a_E>0\). Thus both
schemes, for all seven frames each, preserve \(n\) and
strictly decrease \(w\).

To lift this comparison through any one-hole context, inspect
each constructor recurrence. Equal child \(n\)-values give
equal parent \(n\)-values, including rows that reset \(n\)
to one. With all \(n\)-values fixed, each parent \(w\)
is strictly increasing in each child's \(w\): its relevant
coefficient is one or a positive \(n\)-value. Therefore a
strict decrease in the selected child's \(w\) remains a
strict decrease at its parent. Induction to the root proves
the compatible-closure assertion.
\end{proof}
\end{annexitem}

\begin{annexitem}{annex:p7.nj.vacuous-normalization}{thm:p07-nj-vacuous-normalization}{The full first-order two-phase normalization strategy}
\begin{proof}
For a strongly beta-normalizing finite target term \(V\), let
\(h_\beta(V)\) be the greatest length of a beta-reduction
sequence beginning there. This is a natural number. There
are finitely many redex positions in a finite term, with
one contractum up to alpha-equivalence per position. If
finite reduction lengths were unbounded, at least one of
these finitely many successors would itself have unbounded
remaining lengths. Repeating that choice, for example using
the first such position in a fixed order, would give an
infinite beta reduction. Strong normalization excludes it.
Thus the maximum exists, and \(V\to_\beta^+W\) implies
\(h_\beta(V)>h_\beta(W)\).

The typing part of \cref{thm:p07-nj-cps-simulation} and the
open-term result \cref{thm:p07-stlc-strong-normalization}
make this height available for \(\mathcal T(e(M))\) for
every actual input derivation. Consider
\[
 \mu(M)=\bigl(h_\beta(\mathcal T(e(M))),w(e(M))\bigr)
                  \in\mathbb N\times\mathbb N,
\]
ordered lexicographically with the height coordinate first.
A principal conversion strictly decreases its first coordinate
by the positive simulation; the second coordinate may grow.
A permutation gives alpha-equivalent target terms, so leaves
the first coordinate fixed, and strictly decreases the second
by \cref{thm:p07-nj-permutation-decrease}. These facts hold
at arbitrary proof positions and include every quantified
principal step and every copied branch.

In Phase I, traverse the finite annotated tree in a fixed
order and contract a principal or permutation redex whenever
one exists, performing the prescribed fresh renamings.
The pair \(\mu\) strictly decreases, so this phase is finite.
Each conversion preserves the conclusion and allowed open
assumptions by the preceding preservation results. When the
phase ends no principal or permutation redex remains.

In Phase II, contract any absurdity-absorption redex until
none remains. Each such contraction replaces
\(E[\mathsf{abort}_A(M)]\) by \(\mathsf{abort}_B(M)\).
It deletes at least the following elimination constructor
and can also delete entire auxiliary subproofs. The number
of proof constructors therefore strictly decreases, proving
termination of this phase.

It remains to show that Phase II does not spoil Phase I
normality. The retained subproof \(M\) is unchanged.
At the replacement boundary its new head is
\(\mathsf{abort}\), not an introduction, case or unpacking
head. Hence the surrounding context does not acquire a
principal or permutation redex at that boundary. A new
permutation at the new abort's own argument would require
\(M\) to have been a case or unpacking. But then the old
inner \(\mathsf{abort}_A(M)\) was already a permutation
redex before the contraction, contrary to Phase I normality.
No other new parent--child adjacency is created. This
argument includes the nested-absurdity row, and proves that
Phase II never requires a restart of Phase I.

For a single well-founded record of the strategy, use
\((1,h_\beta(\mathcal T(e(M))),w(e(M)),0)\) in Phase I
and \((0,0,0,s(M))\) in Phase II, where \(s\) counts
proof constructors. Lexicographic order on \(\mathbb N^4\)
decreases at every conversion and at the one phase switch.
The strategy need not compute the beta height; that height
proves termination of the finite-tree traversal rule.

The resulting derivation has none of the three specified
redex classes, so is normal in the declared sense. Every
step retains its formula annotations, conclusion and
eigenvariable conditions and adds no open assumption
formula. Composing these finite preservation statements
gives exactly the theorem's conclusion. The proof concerns
the stated two-phase strategy; it uses no unproved claim
about arbitrary mixtures with absurdity absorption, no
arrow-only substitution for the source logic, and no
quantifier-free surrogate for its actual derivations.
\end{proof}
\end{annexitem}

\begin{annexitem}{annex:p7.nj.principal-permutation-normalization}{thm:p07-nj-principal-permutation-normalization}{The exact first-phase consequence}
\begin{proof}
Apply only Phase I of the preceding construction. Its
lexicographic measure \(\mu\) proves termination without
any absurdity absorption. Its stopping condition is exactly
the absence of principal and permutation redexes, and every
one of its steps preserves the conclusion and allowed open
assumptions. The resulting intermediate derivation therefore
has all the asserted properties. No statement about a
different conversion relation is needed for this consequence.
\end{proof}
\end{annexitem}
